\section{深水区：数学直觉、Lean 生态与研究协作}

前面几章把 AI for Math 拆成了技术栈，但洪乐潼访谈真正有趣的地方在于：她并不把数学看成单纯的“题目集合”。数学研究里有审美、有失败、有长期投入，也有共同体认可。AI 如果只会在已知题库里生成答案，距离数学研究仍然很远；如果能帮助数学家发现结构、形式化证明、补全缺失 lemma、检验反例，它才真正进入研究过程。

\subsection{为什么数学不是普通 NLP 任务}

普通 NLP 任务常常允许语义近似：摘要可以有不同措辞，翻译可以有多个等价表达，聊天回答也允许风格差异。数学证明不同：一个关键条件漏掉，整个证明就可能失败。形式系统把这种严格性推到极致，要求每一个对象、类型、依赖和推理步骤都合法。

\begin{importantbox}{数学任务的特殊性}
数学同时拥有强形式约束和强创造性。强形式约束来自 proof assistant 与定理验证；强创造性来自猜想、审美、定义选择和证明路线设计。AI for Math 必须同时处理这两端。
\end{importantbox}

\subsection{Lean 生态为什么重要}

Lean 的价值不只是验证器，而是生态：mathlib 定理库、社区贡献、形式化风格、tactic 传统、教学材料和数学家参与。模型若要写 Lean 证明，需要理解库里有什么、定理如何命名、哪类 tactic 常用、失败目标如何解读。这和普通代码补全很像，但约束更严格。

\begin{knowledgebox}{Lean 与代码的相似和不同}
相似处在于二者都有语法、类型、库、错误信息和可执行反馈；不同处在于 Lean 的目标不是运行程序，而是证明命题。程序可以靠测试覆盖常见情况，证明则必须满足所有情况。
\end{knowledgebox}

\subsection{Proof Search 的搜索空间}

证明搜索不是线性生成。一个定理可能需要先构造中间 lemma，选择归纳变量，转换目标形式，调用已有定理，或先证明反例不存在。搜索空间巨大，错误路径也很多。AI 的优势在于可以并行尝试和利用语义先验；劣势在于容易生成看似合理但库中不可用的步骤。

\begin{warningbox}{证明搜索的常见失败}
模型可能知道数学思路，却不知道 Lean 库里对应定理叫什么；也可能知道 tactic，却把目标变成更难的形式；还可能生成自然语言上成立但形式系统缺少前提的证明。失败信息本身就是训练数据。
\end{warningbox}

\subsection{本章小结}

AI for Math 的核心不是“模型更聪明一点”，而是把数学研究拆成可反馈的闭环：直觉提出方向，形式化约束过程，验证器提供硬反馈，数学家判断意义。

